% SOURCE_STATEMENT: "Develop a general method for constructing rational and trigonometric $K$-matrices satisfying the matrix reflection equations associated with twisted Yangians and quantum-symmetric-pair coideal subalgebras."
% Preparation and final self-audit: CONTRIBUTING.md prompts C and D; source/audit artifacts retained in the local physics_statements_2026_10_11 batch.
% !TeX program = lualatex
\documentclass[11pt]{article}
\usepackage[margin=25mm]{geometry}
\usepackage{fontspec}
\setmainfont{Latin Modern Roman}
\usepackage{amsmath,amssymb,mathtools}
\usepackage{unicode-math}
\setmathfont{Latin Modern Math}
\usepackage{xurl,hyperref}
\hypersetup{colorlinks=true,urlcolor=blue}
\setcounter{secnumdepth}{0}
\setlength{\parindent}{0pt}
\setlength{\parskip}{5pt}
\setlength{\emergencystretch}{3em}
\newcommand{\sref}[2]{\hyperlink{source:#1:#2}{[S#2]}}
\newcommand{\cataloguescope}[1]{\paragraph{Recorded target.}#1}
\newcommand{\R}{\mathbb R}
\newcommand{\C}{\mathbb C}
\newcommand{\N}{\mathbb N}
\newcommand{\Z}{\mathbb Z}
\newcommand{\Q}{\mathbb Q}
\newcommand{\F}{\mathbb F}
\newcommand{\eps}{\varepsilon}
\begin{document}
\section{Construction of K-Matrices for Coideal Subalgebras}\label{statement}

\subsection{Definitions and mathematical statement}

\paragraph{Shared notation.}$\mathbb N=\{1,2,\ldots\}$; the standard number fields have their usual meanings. All additional conventions are stated locally.


\paragraph{Definitions and conventions.}
An $R$-matrix encodes the exchange relation of an integrable model. A boundary $K$-matrix solves the corresponding reflection equation; one common additive-parameter convention is $R_{12}(u-v)K_1(u)R_{21}(u+v)K_2(v)=K_2(v)R_{12}(u+v)K_1(u)R_{21}(u-v)$. Subscripts mark tensor legs and $R_{21}=PR_{12}P$. A coideal subalgebra $B$ of a Hopf algebra $A$ has coproduct taking values in $B\otimes A$ or $A\otimes B$, according to convention.

\paragraph{Target retained from the canonical source catalogue.}
\paragraph{Catalogue wording retained for traceability.}
Develop a general method for constructing rational and trigonometric $K$-matrices satisfying the matrix reflection equations associated with twisted Yangians and quantum-symmetric-pair coideal subalgebras.

\paragraph{Source scope and qualification.}
The displayed equation only explains an ordinary reflection-equation convention. Twists, representations and multiplicative spectral parameters must be those of the report; both twisted Yangian and quantum-symmetric-pair cases remain in the target.
\subsection{Short English statement}
Find a general way to construct the rational and trigonometric boundary matrices for the indicated coideal algebras.
\subsection{Sources}
\begin{itemize}
\item[\hypertarget{source:statement:1}{S1}] ULAM.AI and the UnsolvedMath Contributors, UnsolvedMath, catalogue entry. Catalogue curation under CC BY 4.0; definitions and English exposition added locally.
\url{https://huggingface.co/datasets/ulamai/UnsolvedMath}.
\item[\hypertarget{source:statement:2}{S2}] Mini-Workshop: Coideal Subalgebras of Quantum Groups (2015), Oberwolfach Reports, DOI 10.4171/owr/2015/10, pp. 533--569. Primary report accessed 2026-10-11; relevant PDF pages 24.
\url{https://doi.org/10.4171/owr/2015/10}.
\item[\hypertarget{source:statement:3}{S3}] Regelskis and Vlaar, reflection matrices (catalogue-listed context; not independently consulted in this preparation).
\url{https://arxiv.org/abs/1602.08471}.
\item[\hypertarget{source:statement:4}{S4}] Recent K-matrix classification (catalogue-listed context; not independently consulted in this preparation).
\url{https://doi.org/10.1016/j.jpaa.2026.108210}.
\end{itemize}
\label{end:statement}
\end{document}
